\documentclass[11pt,a4paper]{article}

\usepackage{amsmath, amssymb, amsthm}
\usepackage{algorithm}
\usepackage{algpseudocode}
\usepackage[margin=1in]{geometry}
\PassOptionsToPackage{hyphens}{url}
\usepackage{hyperref}
\usepackage{listings}
\usepackage{xcolor}
\usepackage{booktabs}
\usepackage{array}
\usepackage{enumitem}
\usepackage{tabularx}
\usepackage{textcomp}

\hypersetup{
  colorlinks=true,
  linkcolor=blue,
  urlcolor=blue,
  citecolor=blue
}

\lstset{
  basicstyle=\ttfamily\footnotesize,
  showstringspaces=false,
  breaklines=true,
  numbers=none,
  frame=single,
  framerule=0.4pt,
  rulecolor=\color{black!50},
  backgroundcolor=\color{black!3},
  columns=fullflexible,
  keepspaces=true,
  upquote=true
}

\lstdefinestyle{python}{
  language=Python,
  keywordstyle=\color{blue!70!black}\bfseries,
  stringstyle=\color{green!50!black},
  commentstyle=\color{gray}\itshape
}

\newtheorem*{definition*}{Definition}
\newtheorem*{remark*}{Remark}

\title{Cost-Sensitive iSFS via Block Coordinate Descent}
\author{Working Notes}
\date{}

\begin{document}
\maketitle

\begin{quote}
A cost-sensitive extension of the \textbf{incomplete Source-Feature Selection (iSFS)} model of Xiang, Yuan, Fan, Wang, Thompson \& Ye (KDD 2013). We introduce a learnable per-class cost vector $C$ and derive the BCD updates for jointly optimizing $(\alpha, \beta, C)$ using only the paper's notation.

\medskip
\noindent\textbf{Sources used throughout:}
\begin{itemize}[leftmargin=*,itemsep=2pt]
  \item \textbf{Xiang et al. (2013).} \emph{Multi-Source Learning with Block-wise Missing Data for Alzheimer's Disease Prediction.} KDD '13. The iSFS objective is eq (14); its data-fitting term is eq (15); the alternating block updates for $\alpha$ and $\beta$ are eq (16) and eqs (17)--(18). Algorithm 2 in \S 3.2 lays out the original BCD scheme.
  \item \textbf{Aggarwal,} \emph{Linear Algebra and Optimization for Machine Learning,} \S 4.10 and \S 8.3. The block coordinate descent framework, the multi-convex monotonicity theorem, and the K-means / ALS examples we inherit.
  \item \textbf{Bertsimas, Pawlowski, Zhuo (JMLR 2018),} \emph{From Predictive Methods to Missing Data Imputation}. Uses BCD vs CD as the inner solver for non-convex imputation.
  \item \textbf{Namkoong \& Duchi (2017),} \emph{Variance-based regularization with convex objectives,} NeurIPS --- distributionally robust optimization with the chi-squared ball, the closed-form max via KKT.
  \item \textbf{Hu, Niu, Sato, Sugiyama (2018),} \emph{Does distributionally robust supervised learning give robust classifiers?,} ICML --- connects DRO directly to cost-sensitive learning.
  \item \textbf{Yuan, Y. (2015),} \emph{Recent advances in trust-region algorithms,} Math.\ Program.\ Ser.\ B, 151, 249--281. Standard formulation of the trust-region subproblem (\S 2, eq.\ 2.3--2.7), Cauchy decrease, predicted / actual reduction ratio, radius update, and bound-constrained affine-scaled (Coleman--Li) trust region (\S 3, eq.\ 3.34--3.37).
\end{itemize}
\end{quote}

\newpage
\tableofcontents
\newpage

%======================================================================
\section{Recap of the iSFS notation we inherit from the paper}
%======================================================================

We use the \textbf{exact symbols} from Xiang et al.\ (2013) \S 3.1. No new variables until \S 3.

\paragraph{Variables.}
\begin{itemize}[leftmargin=*,itemsep=2pt]
  \item $S$ --- number of data sources (e.g., PET, MRI, CSF; so $S = 3$).
  \item $m$ --- a profile index. A profile is the binary indicator of which sources are present, encoded as a decimal integer.
  \item $\mathbf{pf}$ --- vector of profiles for all participants. $|\mathbf{pf}|$ is the number of distinct profiles.
  \item $n_m$ --- number of samples in profile $m$.
  \item $n$ --- number of rows of $\mathbf{X}$ in a generic context (the paper uses both $n$ and $n_m$).
  \item $\mathbf{X}^i \in \mathbb{R}^{n \times p_i}$ --- data matrix for source $i$.
  \item $\mathbf{X}^i_m$ --- data matrix $\mathbf{X}^i$ restricted to samples in profile $m$.
  \item $\boldsymbol{\beta}^i \in \mathbb{R}^{p_i}$ --- feature coefficient block for source $i$. \textbf{Shared across profiles.}
  \item $\boldsymbol{\beta} = (\boldsymbol{\beta}^1, \ldots, \boldsymbol{\beta}^S) \in \mathbb{R}^{p_1 + \ldots + p_S}$ --- full feature coefficient vector.
  \item $\alpha^i_m$ --- scalar weight on source $i$ for profile $m$. \textbf{Profile-specific.}
  \item $\boldsymbol{\alpha}_m = (\alpha^1_m, \ldots, \alpha^S_m) \in \mathbb{R}^S$ --- full source-weight vector for profile $m$.
  \item $\mathbf{y}_m$ --- label vector for samples in profile $m$.
  \item $\mathbf{L}(\cdot, \cdot)$ --- convex loss function (least-squares, logistic, etc.). The paper says: \emph{``L can be any convex loss function such as the least squares loss function or the logistic loss function and $n$ is number of rows of $\mathbf{X}$.''}
  \item $\mathbf{R}_{\boldsymbol{\alpha}}, \mathbf{R}_{\boldsymbol{\beta}}$ --- regularizers on $\boldsymbol{\alpha}$ and $\boldsymbol{\beta}$.
  \item $\lambda$ --- regularization strength on $\mathbf{R}_{\boldsymbol{\beta}}$.
\end{itemize}

\paragraph{Equation (14)} --- the iSFS objective:
\begin{equation}
\min_{\boldsymbol{\alpha},\,\boldsymbol{\beta}}\;\;\frac{1}{|\mathbf{pf}|}\sum_{m \in \mathbf{pf}} f(\mathbf{X}_m,\boldsymbol{\beta},\boldsymbol{\alpha}_m,\mathbf{y}_m) \;+\; \lambda\,\mathbf{R}_{\boldsymbol{\beta}}(\boldsymbol{\beta})
\end{equation}
\[
\text{s.t.}\quad \mathbf{R}_{\boldsymbol{\alpha}}(\boldsymbol{\alpha}_m) \leq 1\;\;\forall\,m \in \mathbf{pf}
\]

\paragraph{Equation (15)} --- the per-profile data-fitting term:
\begin{equation}
f(\mathbf{X},\boldsymbol{\beta},\boldsymbol{\alpha},\mathbf{y}) \;=\; \frac{1}{n}\,\mathbf{L}\!\left(\sum_{i=1}^{S}\alpha^{i}\mathbf{X}^{i}\boldsymbol{\beta}^{i},\;\mathbf{y}\right)
\end{equation}

\paragraph{Equation (16)} --- the $\boldsymbol{\alpha}$ subproblem (when $\boldsymbol{\beta}$ is fixed). Decouples by profile:
\begin{equation}
\min_{\boldsymbol{\alpha}}\;\;\Big\|\sum_{i=1}^{S}\alpha^{i}\mathbf{X}^{i}\boldsymbol{\beta}^{i} - \mathbf{y}\Big\|_{2}^{2}
\quad\text{s.t.}\quad \mathbf{R}_{\boldsymbol{\alpha}}(\boldsymbol{\alpha}) \leq 1
\end{equation}

\paragraph{Equation (17)} --- the $\boldsymbol{\beta}$ subproblem (when $\boldsymbol{\alpha}$ is fixed):
\begin{equation}
\min_{\boldsymbol{\beta}}\;\;g(\boldsymbol{\beta}) + \lambda\,\mathbf{R}_{\boldsymbol{\beta}}(\boldsymbol{\beta}),
\qquad
g(\boldsymbol{\beta}) \;=\; \frac{1}{|\mathbf{pf}|}\sum_{m \in \mathbf{pf}}\frac{1}{2 n_m}\Big\|\sum_{i=1}^{S}\alpha^{i}_{m}\mathbf{X}^{i}_{m}\boldsymbol{\beta}^{i} - \mathbf{y}_m\Big\|_{2}^{2}
\end{equation}

\paragraph{Equation (18)} --- the gradient of $g$ for the $\boldsymbol{\beta}$ step (used by FISTA):
\begin{equation}
\nabla g(\boldsymbol{\beta}^{i}) \;=\; \frac{1}{|\mathbf{pf}|}\sum_{m \in \mathbf{pf}}\frac{1}{n_m}\,\mathbf{I}\!\left(m\,\&\,2^{S-i}\neq 0\right)\,(\alpha^{i}_{m}\mathbf{X}^{i}_{m})^{T}\!\Big(\sum_{j=1}^{S}\alpha^{j}_{m}\mathbf{X}^{j}_{m}\boldsymbol{\beta}^{j} - \mathbf{y}_m\Big)
\end{equation}

The bitwise indicator $\mathbf{I}(m \,\&\, 2^{S-i} \neq 0)$ equals $1$ iff source $i$ is present in profile $m$. This is a literal bit-test on the profile integer.

\textbf{Algorithm 2 in the paper} is the BCD that alternates eq (16) and eq (17) until the objective stops decreasing. It is exactly the multi-convex BCD pattern from Aggarwal \S 4.10.2 --- convex in each block given the other fixed, monotonic decrease across the alternation.

That is the entire toolkit we inherit. Everything below adds \textbf{one new variable $C$} and derives the corresponding new block update.

%======================================================================
\section{Why a cost-sensitive extension --- and why fixed $C$ is not enough}
%======================================================================

\subsection{Motivation}

In Alzheimer's disease prediction (and most clinical settings), the minority class (e.g., progressive MCI converters) is rare. Standard iSFS minimizes average loss and quietly becomes majority-biased. The classical fix is \textbf{cost-sensitive learning}: assign higher misclassification penalty to minority-class errors. From Elkan (2001):
\[
\text{Total Cost} \;=\; C(0,1)\cdot\text{FN} \;+\; C(1,0)\cdot\text{FP},\qquad C(0,1) \gg C(1,0).
\]

For our setting, define a \textbf{per-class cost vector} $C = (C_1, \ldots, C_K)$ where $K$ is the number of classes ($K = 2$ for binary). Each sample's loss contribution is multiplied by $C_{y_j}$ --- the cost weight that corresponds to its class label.

\subsection{Fixed $C$ (classical Elkan-style) as a reference point}

\paragraph{The fixed-$C$ rule.} The standard ``inverse class frequency'' choice:
\[
C_k^{\,0} \;=\; \frac{N}{K \cdot n_k},\qquad k = 1, \ldots, K
\]
where:
\begin{itemize}[leftmargin=*,itemsep=2pt]
  \item $N = \sum_k n_k$ is the total sample count.
  \item $n_k$ is the number of samples in class $k$.
\end{itemize}

This makes minority-class samples count more in the loss. For our running example with $n_- = 12$, $n_+ = 4$, $N = 16$, $K = 2$:
\[
C_-^{\,0} = \frac{16}{2\cdot 12} \approx 0.667,\qquad C_+^{\,0} = \frac{16}{2\cdot 4} = 2.0.
\]
The minority class gets $n_-/n_+ = 3\times$ the per-sample weight of the majority class.

\paragraph{The cost-sensitive version of $\mathbf{L}$.} Promote the paper's loss function $\mathbf{L}$ to its class-weighted form:
\[
\mathbf{L}^{C}(\hat{\mathbf{y}},\mathbf{y}) \;=\; \frac{1}{2}\sum_{j=1}^{n} C_{y_{j}}\,(\hat{y}_{j} - y_{j})^{2}
\]

Plug this into the iSFS objective in place of $\mathbf{L}$. With $C$ fixed at $C^0$, the entire optimization remains in $(\boldsymbol{\alpha}, \boldsymbol{\beta})$:
\[
\min_{\boldsymbol{\alpha},\,\boldsymbol{\beta}}\;\;\frac{1}{|\mathbf{pf}|}\sum_{m \in \mathbf{pf}} f^{C^{0}}(\mathbf{X}_m,\boldsymbol{\beta},\boldsymbol{\alpha}_m,\mathbf{y}_m) \;+\; \lambda\,\mathbf{R}_{\boldsymbol{\beta}}(\boldsymbol{\beta})
\]
where
\[
f^{C}(\mathbf{X},\boldsymbol{\beta},\boldsymbol{\alpha},\mathbf{y}) \;=\; \frac{1}{n}\,\mathbf{L}^{C}\!\left(\sum_{i=1}^{S}\alpha^{i}\mathbf{X}^{i}\boldsymbol{\beta}^{i},\;\mathbf{y}\right).
\]

\textbf{Algorithm 2 from the paper runs unchanged} --- only the inside of $\mathbf{L}$ carries the $C^0$ weights.

\subsection{Why fixed $C$ is not enough --- and what a bounded set buys us}

Define the per-class loss contribution at the current $(\boldsymbol{\alpha}, \boldsymbol{\beta})$:
\[
L_k(\boldsymbol{\alpha},\boldsymbol{\beta}) \;=\; \frac{1}{|\mathbf{pf}|}\sum_{m \in \mathbf{pf}}\frac{1}{2 n_m}\!\!\sum_{\substack{j: y_{m,j} = k}}\!(\hat{y}_{m,j} - y_{m,j})^{2}\;\ge\;0.
\]

The data-fitting part of the objective as a function of $C$ alone is
\[
\Phi_C(C) \;=\; \sum_{k=1}^{K} C_k\,L_k(\boldsymbol{\alpha},\boldsymbol{\beta}).
\]

This is \textbf{linear in $C$} with non-negative coefficients. Unconstrained, $\max_C \Phi_C$ is $+\infty$ and $\min_C \Phi_C$ is $-\infty$ (or $0$ if $C_k \ge 0$), so a ``fix $\boldsymbol{\alpha},\boldsymbol{\beta}$, solve $C$'' step is degenerate without a bounded feasible set. Classical cost-sensitive learning sidesteps the issue by fixing $C$ at a prior (Elkan 2001, Domingos 1999, King \& Zeng 2001, scikit-learn's \texttt{class\_weight='balanced'}).

The DRO observation (Hu et al.\ 2018; Namkoong \& Duchi 2017) is that \textbf{any bounded neighbourhood of $C^0$ in a suitable norm removes the degeneracy}: linear-on-a-compact-convex set has a well-defined maximizer. The simplest such neighbourhood is a ball
\[
\mathcal B(C^0,\,\Delta) \;=\; \Big\{ C \in \mathbb{R}^K \;:\; \tfrac12\|C - C^0\|_2^2 \le \Delta \Big\}.
\]

Once we are inside a ball, we are in the native setting of \textbf{trust-region methods} (Yuan 2015, \S 2): the same ball that bounds the problem also governs an iterative step-length strategy. \S 4 uses this equivalence to turn the $C$-block into a proper TR iteration.

%======================================================================
\section{The cost-sensitive iSFS objective}
%======================================================================

\textbf{One new variable} beyond what the paper uses:
\begin{itemize}[leftmargin=*,itemsep=2pt]
  \item $C \in \mathbb{R}^K_{\ge 0}$ --- per-class cost vector. $K$ parameters total (just $2$ for binary classification).
  \item $C^0 \in \mathbb{R}^K_{\ge 0}$ --- baseline (prior) cost vector. We use the rescaled inverse-frequency baseline (defined in \S 4).
  \item $\Delta_{\max} > 0$ --- \textbf{maximum trust-region radius}: the squared-distance budget $\tfrac12\|C - C^0\|_2^2 \le \Delta_{\max}$ (equivalently, $\|C - C^0\|_2 \le \sqrt{2\Delta_{\max}}$). The TR iteration (\S 7) starts with a smaller working radius $\Delta_0 \le \Delta_{\max}$ and adapts it.
  \item $\mathcal B(C^0,\Delta_{\max}) \;=\; \{C : \tfrac12\|C-C^0\|_2^2 \le \Delta_{\max}\}$ --- problem-level feasible ball for $C$.
\end{itemize}

\noindent\emph{Notation.} We use $\Delta$ (with indices) throughout for trust-region radii, following Yuan (2015, \S 2), and $r_k$ for the actual/predicted reduction ratio in the TR radius-update rule (Yuan denotes it $\rho_k$; we reserve $\rho$ for other uses).

\paragraph{Cost-sensitive loss function} --- extends the paper's $\mathbf{L}$ by attaching the class-cost weight to each sample's residual:
\begin{equation}
\boxed{\;\mathbf{L}^{C}(\hat{\mathbf{y}},\mathbf{y}) \;=\; \frac{1}{2}\sum_{j=1}^{n} C_{y_{j}}\,(\hat{y}_{j} - y_{j})^{2}\;}
\end{equation}

Read this in three pieces:
\begin{itemize}[leftmargin=*,itemsep=2pt]
  \item $(\hat{y}_j - y_j)^2$ --- squared residual for sample $j$ (same as paper's least-squares).
  \item $C_{y_j}$ --- look up the cost weight for sample $j$'s class label $y_j$. This is K-way indexing, \textbf{not} a per-sample free parameter. $C$ still has only $K$ total parameters.
  \item $\sum_j$ --- sum over all samples.
\end{itemize}

When $C_k = 1$ for all $k$, $\mathbf{L}^C$ reduces to $\frac{1}{2}\|\hat{\mathbf{y}} - \mathbf{y}\|^2$, which is exactly the paper's least-squares $\mathbf{L}$.

\paragraph{Cost-sensitive data-fitting term} --- extends eq (15) by replacing $\mathbf{L}$ with $\mathbf{L}^C$:
\begin{equation}
\boxed{\;f^{C}(\mathbf{X},\boldsymbol{\beta},\boldsymbol{\alpha},\mathbf{y}) \;=\; \frac{1}{n}\,\mathbf{L}^{C}\!\left(\sum_{i=1}^{S}\alpha^{i}\mathbf{X}^{i}\boldsymbol{\beta}^{i},\;\mathbf{y}\right)\;}
\end{equation}

\paragraph{Cost-sensitive iSFS objective} --- saddle-point extension of eq (14):
\begin{equation}
\boxed{\;
\min_{\boldsymbol{\alpha},\,\boldsymbol{\beta}}\;\max_{C \in \mathcal B(C^0,\Delta_{\max}),\;C\ge 0}\;\;\frac{1}{|\mathbf{pf}|}\sum_{m \in \mathbf{pf}} f^{C}(\mathbf{X}_m,\boldsymbol{\beta},\boldsymbol{\alpha}_m,\mathbf{y}_m) \;+\; \lambda\,\mathbf{R}_{\boldsymbol{\beta}}(\boldsymbol{\beta})
\;}
\end{equation}
\[
\text{s.t.}\quad \mathbf{R}_{\boldsymbol{\alpha}}(\boldsymbol{\alpha}_m) \leq 1\;\;\forall\,m \in \mathbf{pf}
\]

\paragraph{Three things to read off:}
\begin{enumerate}[leftmargin=*,itemsep=2pt]
  \item \textbf{$\min$ over $(\boldsymbol{\alpha}, \boldsymbol{\beta})$, $\max$ over $C$.} The saddle-point structure used in DRO; $(\boldsymbol{\alpha},\boldsymbol{\beta})$ is trained against the worst-case class re-weighting within the ball.
  \item \textbf{$C$ lives in a ball around $C^0$.} $C$ is restricted to be within squared-distance $\Delta_{\max}$ of the prior $C^0$ and non-negative. The ball alone is enough to bound $\Phi_C$ because the objective is linear in $C$.
  \item \textbf{The only changes from eq (14):} new variable $C$; $f \to f^C$; new constraint $C \in \mathcal B(C^0,\Delta_{\max}) \cap \mathbb{R}^K_{\ge 0}$. Everything else is identical.
\end{enumerate}

\textbf{Why a max over $C$, not a min.} A pure-min framing in $C$ assigns less weight to the harder class, the opposite of cost-sensitive intent. The max framing forces the model to do well against whichever class the data currently makes hard.

%======================================================================
\section{The feasible set for $C$: a pure trust-region ball}
%======================================================================

The feasible set for $C$ in the problem statement (\S 3) is a single ball around the prior $C^0$:
\[
\mathcal B(C^0, \Delta_{\max}) \cap \mathbb{R}^K_{\ge 0}
\;=\; \Big\{\,C \in \mathbb{R}^K_{\ge 0} \;:\; \tfrac12\|C - C^0\|_2^2 \leq \Delta_{\max}\,\Big\}.
\]

\subsection{Why the ball alone is enough}

A ball of finite squared-radius $\Delta_{\max}$ around a finite center $C^0$ bounds every component of $C$, so $\max_{C \in \mathcal B(C^0,\Delta_{\max})} L^\top C$ is attained at a unique point on the boundary (Cauchy--Schwarz). The ball is the problem-level feasible set, and the same ball is the native geometric object for iterative trust-region optimization (Yuan 2015, \S 2). One ball plays two roles: outer feasibility in $\max_C$, inner trust region per TR step.

\subsection{The baseline $C^0$ (rescaled inverse frequency)}

Same prior as before:
\[
C_k^{\,0} \;=\; \frac{K \cdot (1/n_k)}{\displaystyle\sum_{j=1}^{K}(1/n_j)}.
\]
For $n_- = 12,\ n_+ = 4,\ K = 2$: $C^0 = (0.5,\,1.5)$. The rescaling is no longer forced by any constraint, but we keep it so the $\Delta_{\max} \to 0$ limit recovers the standard balanced baseline used by \texttt{class\_weight=`balanced'}.

\subsection{Problem radius $\Delta_{\max}$ vs.\ working radius $\Delta_k$}

Two distinct quantities, both in squared-distance units:

\begin{itemize}[leftmargin=*,itemsep=2pt]
  \item $\Delta_{\max}$ --- \textbf{problem-level} TR radius appearing in the saddle-point objective (\S 3). A hyperparameter; the maximum deviation from $C^0$ ever allowed.
  \item $\Delta_k$ --- \textbf{iteration} TR radius at step $k$ of the $C$-block. Starts at some $\Delta_0 \le \Delta_{\max}$, is adapted by the Yuan (2015) actual/predicted ratio rule, and is clipped to $[\Delta_{\min},\,\Delta_{\max}]$.
\end{itemize}

$\Delta_{\max}$ controls \emph{how far} the adversarial $C$ is allowed to move. $\Delta_k$ controls \emph{how big each step} is inside that allowance. This is the standard separation in trust-region methods (Yuan 2015, \S 2).

\subsection{What the radius means in plain terms}

\begin{center}
\begin{tabular}{ll}
\toprule
\textbf{$\Delta_{\max}$ value} & \textbf{behavior of $C$} \\
\midrule
$\Delta_{\max} \to 0$ & $C$ is pinned to $C^0$ --- recovers classical fixed cost-sensitive learning. \\
$\Delta_{\max}$ small & $C$ stays close to $C^0$: ``almost-classical'' adaptation. \\
$\Delta_{\max}$ moderate & $C$ adapts to per-class loss but is anchored to the prior. \\
$\Delta_{\max}$ large & $C$ is free to roam up to the non-negative-orthant boundary. \\
\bottomrule
\end{tabular}
\end{center}

Tune $\Delta_{\max}$ by cross-validation. $\Delta_0$ is typically set to $\Delta_{\max}$ or a modest fraction of it; the Yuan rule shrinks it automatically if the linear model over-predicts reduction.

\subsection{Where this comes from --- research}

Two orthogonal streams inform this construction:

\begin{itemize}[leftmargin=*,itemsep=2pt]
  \item \textbf{Distributionally robust optimization (DRO).} Hu, Niu, Sato, Sugiyama (2018, ICML) and Namkoong \& Duchi (2017, NeurIPS) justify an $\ell_2$ / $\chi^2$ ball on a per-class re-weighting vector as the principled way to \emph{learn} cost weights instead of fixing them. The ball is the constraint set; the ``adversary'' inside the $\max$ picks the worst re-weighting within it.
  \item \textbf{Trust-region methods.} Yuan (2015, Math.\ Program.) surveys the canonical TR recipe: build a local model $m_k(d)$ around the current iterate, trust it only within a ball of radius $\Delta_k$, accept or reject the step via the ratio of actual to predicted reduction, and adapt $\Delta_k$ (eq.\ 2.3--2.7, p.\ 252--253). \S 3 of the survey extends this to bound-constrained problems via the Coleman--Li affine-scaled ellipsoidal TR (eq.\ 3.34--3.37, p.\ 262), which is how we handle $C \ge 0$.
\end{itemize}

\subsection{Summary}

Problem level: $C$ lives in a single ball around $C^0$ (plus $C \ge 0$). Algorithmically, $C$ is updated by a trust-region iteration whose radius $\Delta_k$ starts at $\Delta_0 \le \Delta_{\max}$ and adapts step-by-step. \textbf{Before} applying TR to $C$, the next section (\S 5) presents the general trust-region framework as it appears in Yuan (2015), in standard optimization notation, so that the reader can verify the machinery independently of our cost-sensitive application.

%======================================================================
\section{Trust-region optimization: the general framework (Yuan 2015)}
\label{sec:tr-general}
%======================================================================

This section reproduces the core trust-region (TR) algorithm from Yuan, Y.\ (2015), \emph{Recent advances in trust-region algorithms}, Math.\ Program.\ Ser.\ B 151, 249--281. Everything here is in \textbf{standard optimization notation} --- minimizing a smooth function $f : \mathbb{R}^n \to \mathbb{R}$, with no reference to cost-sensitive learning or iSFS. The purpose is to establish the equations we will use in \S 7 when we apply TR to the $C$-block.

\subsection{The problem and the idea}

We want to solve
\[
\min_{x \in \mathbb{R}^n}\; f(x)
\]
where $f$ is smooth (twice continuously differentiable). At the current iterate $x_k$ we know $f(x_k)$, the gradient $g_k = \nabla f(x_k)$, and optionally a Hessian or Hessian-surrogate $B_k \approx \nabla^2 f(x_k)$.

The \textbf{trust-region idea} (Yuan \S 2, p.\ 251--253): instead of taking a full Newton or gradient step, build a \emph{local model} $m_k(d)$ of $f(x_k + d)$ around $d = 0$, trust it only within a ball of radius $\Delta_k$, and solve:
\begin{equation}
\boxed{\;
\min_{d \in \mathbb{R}^n}\; m_k(d) \;=\; f(x_k) + g_k^\top d + \tfrac12\,d^\top B_k\,d
\qquad\text{s.t.}\quad \|d\|_2 \;\le\; \Delta_k.
\;}
\tag{Yuan, eq.\ 2.3--2.4}
\end{equation}

\begin{itemize}[leftmargin=*,itemsep=2pt]
  \item $m_k(d)$ is a quadratic model that matches $f$ at $d = 0$: $m_k(0) = f(x_k)$, $\nabla m_k(0) = g_k$.
  \item $B_k$ can be the exact Hessian $\nabla^2 f(x_k)$, a quasi-Newton approximation (BFGS, L-BFGS), or even $B_k = 0$ (giving a \emph{linear model}).
  \item $\Delta_k > 0$ is the \textbf{trust-region radius}: how far we allow $d$ to go. It is adapted at every iteration based on how well the model predicted the actual function value.
\end{itemize}

\subsection{Predicted reduction, actual reduction, and the ratio $r_k$}

After solving the subproblem to get $d_k$, we compare what the model predicted to what actually happened (Yuan eq.\ 2.7, p.\ 253):
\begin{equation}
\boxed{\;
\text{Pred}_k \;=\; m_k(0) - m_k(d_k),\qquad
\text{Ared}_k \;=\; f(x_k) - f(x_k + d_k),\qquad
r_k \;=\; \frac{\text{Ared}_k}{\text{Pred}_k}.
\;}
\tag{Yuan, eq.\ 2.7}
\end{equation}

\begin{itemize}[leftmargin=*,itemsep=2pt]
  \item $\text{Pred}_k$ is the reduction the model predicts: always $\ge 0$ if $d_k$ is a descent direction for $m_k$.
  \item $\text{Ared}_k$ is the actual reduction in the true function $f$.
  \item $r_k$ measures agreement between model and reality:
    \begin{itemize}
      \item $r_k \approx 1$: model is accurate. Trust it more $\to$ expand $\Delta_k$.
      \item $r_k$ small or negative: model is inaccurate. Trust it less $\to$ shrink $\Delta_k$.
    \end{itemize}
\end{itemize}

\subsection{Step acceptance and radius update}

\paragraph{Step acceptance.} Accept $x_{k+1} = x_k + d_k$ if $r_k \ge \tau_1$; otherwise reject the step and keep $x_{k+1} = x_k$ (Yuan \S 2, p.\ 253).

\paragraph{Radius update rule.} Adapt $\Delta_k$ based on $r_k$ (Yuan eq.\ 2.8--2.9, p.\ 253):
\begin{equation}
\boxed{\;
\Delta_{k+1} \;=\;
\begin{cases}
\gamma_2\,\Delta_k & \text{if } r_k \ge \tau_2 \quad\text{(model is very accurate: expand)}\\[3pt]
\Delta_k & \text{if } \tau_1 \le r_k < \tau_2 \quad\text{(model is adequate: keep)}\\[3pt]
\gamma_1\,\Delta_k & \text{if } r_k < \tau_1 \quad\text{(model is poor: shrink)}
\end{cases}
\;}
\tag{Yuan, eq.\ 2.8--2.9}
\end{equation}
with standard parameters (used throughout Yuan and in our Algorithm 3):
\[
\tau_1 = 0.25,\qquad \tau_2 = 0.75,\qquad \gamma_1 = 0.5,\qquad \gamma_2 = 2.
\]

\subsection{The Cauchy decrease condition}

For convergence theory, the TR step $d_k$ must achieve at least a fixed fraction of the \textbf{Cauchy point} --- the steepest-descent step clipped to the ball. The sufficient-descent condition is (Yuan eq.\ 2.5, p.\ 252):
\begin{equation}
\boxed{\;
\text{Pred}_k \;\ge\; c\,\|g_k\|_2\,\min\!\left\{\Delta_k,\;\frac{\|g_k\|_2}{\|B_k\|_2}\right\},\qquad c > 0.
\;}
\tag{Yuan, eq.\ 2.5}
\end{equation}

When $B_k = 0$ (our case), $\|B_k\|_2 = 0$, so the $\min$ reduces to $\Delta_k$:
\[
\text{Pred}_k \;\ge\; c\,\|g_k\|_2\,\Delta_k.
\]
Since $\text{Pred}_k = g_k^\top d_k = \|g_k\|_2 \cdot \Delta_k$ for the Cauchy step of a linear model, this condition is always satisfied with $c = 1$.

\subsection{Special case: linear model ($B_k = 0$)}

When the model is purely linear ($B_k = 0$), the quadratic term vanishes:
\[
m_k(d) \;=\; f(x_k) + g_k^\top d.
\]
The TR subproblem becomes
\[
\min_{\|d\|_2 \le \Delta_k}\;\; g_k^\top d.
\]
The minimizer points along $-g_k$ (steepest descent direction) with step length $\Delta_k$:
\begin{equation}
\boxed{\;d_k^\star \;=\; -\,\Delta_k\,\frac{g_k}{\|g_k\|_2}.}
\tag{Cauchy step, linear model}
\end{equation}

Verify: $\|d_k^\star\|_2 = \Delta_k$ (on the boundary). Predicted reduction: $\text{Pred}_k = -g_k^\top d_k^\star = \Delta_k\,\|g_k\|_2 > 0$.

For a \textbf{maximization} problem $\max_x f(x)$ (equivalently $\min_x -f(x)$), the step reverses:
\[
d_k^\star \;=\; +\,\Delta_k\,\frac{g_k}{\|g_k\|_2}
\]
and the ratio uses $r_k = (f(x_k + d_k) - f(x_k))/(m_k(d_k) - m_k(0))$, comparing actual increase to predicted increase.

\paragraph{Why the linear model is not trivial in an alternating scheme.} If $f$ were truly linear, one step to the boundary of $\Delta_k$ would solve the subproblem exactly ($r_k = 1$ always) and there would be no reason to iterate. The value of TR emerges when $f$ is \emph{approximately} linear locally but \emph{nonlinear globally} --- i.e., the gradient changes as we move. In block coordinate descent, the gradient of one block's objective changes whenever the other blocks are re-optimized. This is exactly our setting: $\nabla_C \Phi = L$ changes after each Block A re-solve of $(\boldsymbol\alpha,\boldsymbol\beta)$. The TR framework handles this gracefully by testing each step's prediction quality and adapting the step size.

\subsection{The full TR algorithm (general form)}

\begin{algorithm}[H]
\caption{General Trust-Region Algorithm (Yuan 2015, \S 2)}
\begin{algorithmic}[1]
\Require $f$, starting point $x_0$, initial radius $\Delta_0 > 0$, thresholds $\tau_1 < \tau_2$, multipliers $0 < \gamma_1 < 1 < \gamma_2$
\Ensure $x^\star$ (approximate minimizer)
\For{$k = 0, 1, 2, \ldots$}
    \State Compute gradient $g_k = \nabla f(x_k)$ and (optionally) Hessian-surrogate $B_k$.
    \State Build model $m_k(d) = f(x_k) + g_k^\top d + \tfrac12 d^\top B_k d$.
    \State Solve the TR subproblem: $d_k \gets \arg\min_{\|d\|_2 \le \Delta_k} m_k(d)$.
    \State Compute $\text{Ared}_k = f(x_k) - f(x_k + d_k)$ and $\text{Pred}_k = m_k(0) - m_k(d_k)$.
    \State $r_k \gets \text{Ared}_k / \text{Pred}_k$.
    \If{$r_k \ge \tau_1$}
        \State Accept: $x_{k+1} \gets x_k + d_k$.
    \Else
        \State Reject: $x_{k+1} \gets x_k$.
    \EndIf
    \State Update radius:
    $\Delta_{k+1} \gets \begin{cases}\gamma_2\Delta_k & r_k \ge \tau_2 \\ \Delta_k & \tau_1 \le r_k < \tau_2 \\ \gamma_1\Delta_k & r_k < \tau_1\end{cases}$
    \If{$\|g_k\|_2 < \epsilon$}
        \State \Return $x_{k+1}$.
    \EndIf
\EndFor
\end{algorithmic}
\end{algorithm}

\textbf{Convergence.} Under the Cauchy-decrease condition (eq.\ 2.5), Yuan (2015, Theorem 2.1) proves:
\[
\liminf_{k \to \infty}\;\|g_k\|_2 \;=\; 0.
\]
The iterates converge to a first-order stationary point of $f$. If $B_k$ also captures second-order information, superlinear convergence rates are possible (Yuan \S 2.2).

\subsection{Extension to bound constraints: Coleman--Li scaling (Yuan \S 3.3)}

For
\[
\min_{x}\; f(x) \qquad\text{s.t.}\quad \ell \;\le\; x \;\le\; u
\]
where $\ell, u \in (\mathbb{R} \cup \{-\infty, +\infty\})^n$ are elementwise bounds, Coleman \& Li (1996) replace the Euclidean ball with an \textbf{affine-scaled ellipsoid} that shrinks in directions where the iterate is near a bound (Yuan eq.\ 3.34--3.37, p.\ 262):

\paragraph{Step 1 --- define the scaling vector.} For each component $i$:
\begin{equation}
\boxed{\;v^{(i)}(x) \;=\;
\begin{cases}
x_i - \ell_i & \text{if } (\nabla f(x))_i \ge 0 \text{ and } \ell_i > -\infty \\[3pt]
u_i - x_i & \text{if } (\nabla f(x))_i < 0 \text{ and } u_i < +\infty \\[3pt]
+\infty & \text{otherwise}
\end{cases}
\;}
\tag{Yuan, eq.\ 3.34--3.35}
\end{equation}

Interpretation: $v^{(i)}$ is the distance from $x_i$ to the nearest \emph{active} bound (the one the gradient is pushing toward). If the gradient pushes away from any finite bound, $v^{(i)} = +\infty$ (no scaling needed).

\paragraph{Step 2 --- build the scaling matrix.}
\begin{equation}
D_k \;=\; \mathrm{diag}\!\Big(\sqrt{|v^{(1)}(x_k)|},\;\ldots,\;\sqrt{|v^{(n)}(x_k)|}\Big).
\tag{Yuan, eq.\ 3.36}
\end{equation}

\paragraph{Step 3 --- replace the Euclidean ball with the ellipsoid.}
\begin{equation}
\boxed{\;\|D_k^{-1}\,d\|_2 \;\le\; \Delta_k\qquad\text{replaces}\qquad \|d\|_2 \;\le\; \Delta_k.}
\tag{Yuan, eq.\ 3.37}
\end{equation}

Net effect: as $x_i \to \ell_i$ (or $u_i$), $v^{(i)} \to 0$, so $D_k^{(i,i)} \to 0$, and $|d_i| / D_k^{(i,i)} \to \infty$ unless $d_i \to 0$ too. The iterate is automatically prevented from crossing the bound. No explicit projection or clipping is needed.

The rest of the TR algorithm (model, ratio, radius update) is identical to the unconstrained case --- only the ball shape changes. The Cauchy decrease condition (eq.\ 2.5) also carries over with $\|\cdot\|_{D_k^{-1}}$ in place of $\|\cdot\|_2$.

\subsection{Connection to our $C$-block}

\S 7 applies this framework to the cost-sensitive $C$-block as follows:
\begin{center}
\begin{tabular}{lll}
\toprule
\textbf{general TR} & \textbf{our $C$-block} & \textbf{why} \\
\midrule
$x_k$ & $C^{(k)}$ & current iterate \\
$f(x)$ & $-F(C) = -\min_{\alpha,\beta}\Phi(\alpha,\beta,C)$ & we maximize $F$, so minimize $-F$ \\
$g_k = \nabla f$ & $-L(\alpha_k^\star,\beta_k^\star)$ & gradient of $-F$ by envelope theorem \\
$B_k$ & $0$ & linear model (no Hessian) \\
$\Delta_k$ & working TR radius (squared-distance budget) & adapted by ratio rule \\
$\ell \le x \le u$ & $C \ge 0$ (lower bound $\ell = 0$, no upper bound) & Coleman--Li for non-negativity \\
\bottomrule
\end{tabular}
\end{center}

\noindent\emph{Convention note.} Yuan uses $\|d\|_2 \le \Delta_k$ where $\Delta_k$ is the Euclidean radius. Throughout the rest of this document we use $\tfrac12\|d\|_2^2 \le \Delta_k$, treating $\Delta_k$ as a \emph{squared-distance budget} so it has the same units as $\Delta_{\max}$ in the problem statement (\S 3). The two conventions are related by $\Delta_k^{\mathrm{Yuan}} = \sqrt{2\Delta_k^{\mathrm{ours}}}$. All equations in \S 7 and Algorithm 3 use our squared-distance convention.

%======================================================================
\section{Step 1 --- fix $C$, solve $(\boldsymbol{\alpha}, \boldsymbol{\beta})$}
%======================================================================

\begin{definition*}[FISTA]
\textbf{FISTA} = \emph{Fast Iterative Shrinkage-Thresholding Algorithm} (Beck \& Teboulle, 2009, \emph{SIAM Journal on Imaging Sciences}). It is the standard \textbf{accelerated proximal gradient method} for solving optimization problems of the form
\[
\min_{x}\;\;f(x) + g(x)
\]
where $f$ is convex and smooth (with Lipschitz-continuous gradient $\nabla f$) and $g$ is convex (possibly non-smooth) but has an efficient \textbf{proximal operator}
\[
\mathrm{prox}_{tg}(v) \;=\; \arg\min_{x}\Big\{\,\tfrac{1}{2}\|x - v\|^2 + t\cdot g(x)\,\Big\}.
\]
Each FISTA step takes a gradient step on $f$, applies the proximal operator of $g$, and uses Nesterov-style momentum to accelerate. It converges at rate $O(1/k^2)$ --- strictly better than vanilla gradient descent's $O(1/k)$. The original iSFS Algorithm 2 in Xiang et al.\ (2013) cites Beck \& Teboulle and uses FISTA for both eq (16) (with $f$ = least-squares, $g$ = unit-ball indicator) and eq (17) (with $f$ = $g(\boldsymbol{\beta})$, $g$ = $\lambda \mathbf{R}_{\boldsymbol{\beta}}(\boldsymbol{\beta})$). We use the same algorithm here, modified only by the per-sample cost weight $C_{y_{m,j}}$ that appears inside $f$.
\end{definition*}

Hold $C$ fixed inside $f^C$. The cost-sensitive iSFS problem reduces to a $C$-weighted version of the paper's eq (14):
\[
\min_{\boldsymbol{\alpha},\,\boldsymbol{\beta}}\;\;\frac{1}{|\mathbf{pf}|}\sum_{m \in \mathbf{pf}} f^{C}(\mathbf{X}_m,\boldsymbol{\beta},\boldsymbol{\alpha}_m,\mathbf{y}_m) \;+\; \lambda\,\mathbf{R}_{\boldsymbol{\beta}}(\boldsymbol{\beta})
\]

For least-squares loss the inner expansion is:
\[
\mathbf{L}^{C}\!\left(\sum_{i}\alpha^{i}_{m}\mathbf{X}^{i}_{m}\boldsymbol{\beta}^{i},\;\mathbf{y}_m\right) \;=\; \frac{1}{2}\sum_{j=1}^{n_m} C_{y_{m,j}}\!\left(\hat{y}_{m,j} - y_{m,j}\right)^{2}
\]

\paragraph{One auxiliary definition:}
\[
\hat{y}_{m,j}(\boldsymbol{\alpha}_m,\boldsymbol{\beta}) \;=\; \sum_{i=1}^{S}\alpha^{i}_{m}\,\mathbf{x}^{i}_{m,j}\!\cdot\boldsymbol{\beta}^{i}
\]

Variables:
\begin{itemize}[leftmargin=*,itemsep=2pt]
  \item $j \in \{1, \ldots, n_m\}$ --- sample index inside profile $m$.
  \item $y_{m,j}$ --- the class label of sample $j$ in profile $m$.
  \item $\mathbf{x}^i_{m,j}$ --- the $j$-th row of $\mathbf{X}^i_m$, i.e., the source-$i$ feature vector of sample $j$ in profile $m$.
  \item $\hat{y}_{m,j}$ --- the model prediction for sample $j$ in profile $m$.
\end{itemize}

$\hat{y}_{m,j}$ is just the per-sample form of $\sum_i \alpha^i_m \mathbf{X}^i_m \boldsymbol{\beta}^i$ from eq (15). It is linear in $\boldsymbol{\alpha}_m$ (with $\boldsymbol{\beta}$ fixed) and linear in $\boldsymbol{\beta}$ (with $\boldsymbol{\alpha}_m$ fixed).

The inner block is itself a 2-block alternation, exactly Algorithm 2 of the paper, with one change: every squared residual carries $C_{y_{m,j}}$.

\subsection{Inner step 1a --- $\boldsymbol{\alpha}$ update (per profile)}

For each profile $m$, the cost-sensitive version of paper eq (16):
\begin{equation}
\boxed{\;\min_{\boldsymbol{\alpha}_m}\;\;\frac{1}{2}\sum_{j=1}^{n_m} C_{y_{m,j}}\!\left(\sum_{i=1}^{S}\alpha^{i}_{m}\,\mathbf{x}^{i}_{m,j}\!\cdot\boldsymbol{\beta}^{i} \;-\; y_{m,j}\right)^{\!2}\quad\text{s.t.}\quad \mathbf{R}_{\boldsymbol{\alpha}}(\boldsymbol{\alpha}_m) \leq 1\;}
\end{equation}

\textbf{Sanity check.} When $C_k = 1$ for all $k$, every $C_{y_{m,j}}$ becomes $1$, the inner sum collapses to $\frac{1}{2}\|\sum_i \alpha^i_m \mathbf{X}^i_m \boldsymbol{\beta}^i - \mathbf{y}_m\|^2$, and we recover paper eq (16) verbatim. \checkmark

\paragraph{Gradient (used by FISTA / proximal gradient).} Take the partial with respect to $\alpha^i_m$:
\[
\frac{\partial}{\partial \alpha^{i}_{m}}\!\left[\tfrac{1}{2}\sum_{j=1}^{n_m} C_{y_{m,j}}(\hat{y}_{m,j} - y_{m,j})^{2}\right] \;=\; \sum_{j=1}^{n_m} C_{y_{m,j}}\,(\hat{y}_{m,j} - y_{m,j})\,\big(\mathbf{x}^{i}_{m,j}\!\cdot\boldsymbol{\beta}^{i}\big)
\]

Read this in three pieces:
\begin{itemize}[leftmargin=*,itemsep=2pt]
  \item $(\hat{y}_{m,j} - y_{m,j})$ --- residual for sample $j$.
  \item $C_{y_{m,j}}$ --- class-cost weight for sample $j$ (this is the cost-sensitive entry point).
  \item $(\mathbf{x}^i_{m,j} \cdot \boldsymbol{\beta}^i)$ --- partial of $\hat{y}_{m,j}$ with respect to $\alpha^i_m$.
\end{itemize}

When $C \equiv 1$ the $C$ factors disappear and this is the gradient of paper eq (16). \checkmark

The unit-ball constraint $\mathbf{R}_{\boldsymbol{\alpha}}(\boldsymbol{\alpha}_m) \leq 1$ is handled by projection inside FISTA, exactly as in the paper's discussion of eq (16).

\subsection{Inner step 1b --- $\boldsymbol{\beta}$ update (global)}

Stack profiles. The cost-sensitive version of paper's $g$ (eq 17):
\[
g^{C}(\boldsymbol{\beta}) \;=\; \frac{1}{|\mathbf{pf}|}\sum_{m \in \mathbf{pf}}\frac{1}{2 n_m}\sum_{j=1}^{n_m} C_{y_{m,j}}\!\left(\hat{y}_{m,j} - y_{m,j}\right)^{2}
\]

The $\boldsymbol{\beta}$ subproblem:
\begin{equation}
\boxed{\;\min_{\boldsymbol{\beta}}\;\;g^{C}(\boldsymbol{\beta}) \;+\; \lambda\,\mathbf{R}_{\boldsymbol{\beta}}(\boldsymbol{\beta})\;}
\end{equation}

\paragraph{Gradient with respect to $\boldsymbol{\beta}^i$} --- cost-sensitive version of paper eq (18):
\begin{equation}
\boxed{\;\nabla g^{C}(\boldsymbol{\beta}^{i}) \;=\; \frac{1}{|\mathbf{pf}|}\sum_{m \in \mathbf{pf}}\frac{1}{n_m}\,\mathbf{I}\!\left(m\,\&\,2^{S-i}\neq 0\right)\sum_{j=1}^{n_m} C_{y_{m,j}}\,(\hat{y}_{m,j} - y_{m,j})\,\alpha^{i}_{m}\,\mathbf{x}^{i}_{m,j}\;}
\end{equation}

Variables in this expression:
\begin{itemize}[leftmargin=*,itemsep=2pt]
  \item $\mathbf{I}(m \,\&\, 2^{S-i} \neq 0)$ --- bitwise indicator from eq (18) of the paper. Equals $1$ iff source $i$ is present in profile $m$. \textbf{Identical to the paper.}
  \item $\alpha^i_m \mathbf{x}^i_{m,j}$ --- the source-$i$ contribution from sample $j$ in profile $m$.
  \item $(\hat{y}_{m,j} - y_{m,j})$ --- residual.
  \item $C_{y_{m,j}}$ --- class-cost weight on the residual.
\end{itemize}

Compare to the paper's eq (18):
\[
\nabla g(\boldsymbol{\beta}^{i}) \;=\; \frac{1}{|\mathbf{pf}|}\sum_{m \in \mathbf{pf}}\frac{1}{n_m}\,\mathbf{I}\!\left(m\,\&\,2^{S-i}\neq 0\right)\,(\alpha^{i}_{m}\mathbf{X}^{i}_{m})^{T}\!\Big(\sum_{j}\alpha^{j}_{m}\mathbf{X}^{j}_{m}\boldsymbol{\beta}^{j} - \mathbf{y}_m\Big)
\]

The only difference: each per-sample term carries an extra $C_{y_{m,j}}$ factor. With $C \equiv 1$ the two gradients are identical. \checkmark

The $\boldsymbol{\beta}$ subproblem is solved by accelerated gradient (FISTA), exactly as in Algorithm 2 of the paper. \textbf{No new machinery --- only the gradient picks up the $C_{y_{m,j}}$ weight.}

%======================================================================
\section{Step 2 --- fix $(\boldsymbol{\alpha}, \boldsymbol{\beta})$, solve $C$ by a trust-region iteration}
%======================================================================

The $C$-block is solved by an iterative trust-region (TR) procedure in the sense of Yuan (2015, \S 2). Because the data-fitting term is linear in $C$, a single TR step has a closed-form inner solve; iteration matters at the outer level, where re-solving Block A at each new $C$ makes the reduced objective genuinely nonlinear in $C$.

\subsection{Reduction to a linear function of $C$}

Group every per-sample term by the class label $y_{m,j}$. Define for each $k \in \{1, \ldots, K\}$:
\begin{equation}
\boxed{\;L_k(\boldsymbol{\alpha},\boldsymbol{\beta}) \;=\; \frac{1}{|\mathbf{pf}|}\sum_{m \in \mathbf{pf}}\frac{1}{2 n_m}\!\!\sum_{\substack{j=1,\ldots,n_m \\ y_{m,j} = k}}\!\!(\hat{y}_{m,j} - y_{m,j})^{2}\;}
\end{equation}

$L_k$ is the \textbf{per-class loss coefficient}. With $(\boldsymbol{\alpha}, \boldsymbol{\beta})$ fixed, it is just a positive number --- the total residual-squared contribution from samples of class $k$, weighted by the same $(1/|\mathbf{pf}|)(1/n_m)$ factors that appear in the iSFS objective.

The data-fitting term is then \textbf{linear in $C$}:
\[
\frac{1}{|\mathbf{pf}|}\sum_{m \in \mathbf{pf}}\frac{1}{n_m}\,\mathbf{L}^{C}(\cdots) \;=\; \sum_{k=1}^{K} C_k\,L_k(\boldsymbol{\alpha},\boldsymbol{\beta})
\]

So at any current $C$ the $C$-block becomes (with $(\boldsymbol\alpha,\boldsymbol\beta)$ fixed):
\begin{equation}
\boxed{\;\max_{C \in \mathcal B(C^0,\Delta_{\max}),\;C\ge 0}\;\;\Phi_C(C)\;=\;\sum_{k=1}^{K} C_k\,L_k\;}
\end{equation}

The function is linear in $C$, so its gradient is $\nabla \Phi_C = L \ge 0$.

\subsection{TR-L: linear model with adaptive radius}

We solve this subproblem by the Yuan (2015) trust-region recipe, using a \textbf{linear local model}. At TR iterate $k$ with current cost vector $C^{(k)}$ and working radius $\Delta_k$, define
\begin{equation}
m_k(d) \;=\; \Phi_C(C^{(k)}) \;+\; L_k^\top d,\qquad L_k \;=\; L(\boldsymbol\alpha^{(k)},\boldsymbol\beta^{(k)})
\end{equation}
where $L_k$ is the per-class loss vector evaluated at the \emph{current} $(\boldsymbol\alpha^{(k)},\boldsymbol\beta^{(k)})$. (If we re-solve Block A inside the TR loop, $L_k$ changes with $k$; if we only recompute it at outer BCD boundaries, $L_k \equiv L$ within a block.)

\paragraph{TR subproblem.} Maximize $m_k$ over a $\Delta_k$-ball, respecting the problem-level ball and $C \ge 0$:
\begin{equation}
\boxed{\;d_k^\star \;=\; \arg\max_{d}\;\;L_k^\top d\qquad\text{s.t.}\;\;\tfrac12\|d\|_2^2 \le \Delta_k,\;\;C^{(k)}+d \ge 0,\;\;\tfrac12\|C^{(k)}+d-C^0\|_2^2 \le \Delta_{\max}.\;}
\end{equation}

Without the bound and outer-ball constraints, the closed form is Yuan's Cauchy step for linear models:
\[
d_k^{\mathrm{lin}} \;=\; \sqrt{2\Delta_k}\;\frac{L_k}{\|L_k\|_2}.
\]
With the two additional constraints, we clip $d_k^\star$ to the intersection of (i) the $\Delta_k$-ball, (ii) the non-negative orthant shifted by $C^{(k)}$, and (iii) the $\Delta_{\max}$-ball around $C^0$. For $K = 2$ this intersection is a one-dimensional line segment along $L_k/\|L_k\|_2$ and the clip is one scalar \texttt{max/min}. See \S 7.4 for the active-set projection that handles bounds rigorously (the Coleman--Li scaling route).

\paragraph{Ratio test and radius update (Yuan 2015, eq.\ 2.7).}
\begin{equation}
r_k \;=\; \frac{\Phi_C(C^{(k)} + d_k^\star) - \Phi_C(C^{(k)})}{m_k(d_k^\star) - m_k(0)}.
\end{equation}
For a linear model $m_k$ and a linear $\Phi_C$ (with $(\boldsymbol\alpha,\boldsymbol\beta)$ held fixed within the block), numerator and denominator are both exactly $L_k^\top d_k^\star$, so $r_k = 1$ identically. This is the expected behavior: when the model is exact, every step is accepted and the working radius can grow. The non-trivial use of $r_k$ is in the \textbf{outer} BCD loop (Algorithm 3), where re-solving Block A for the new $C$ changes $L$ --- then the actual reduction (using the re-optimized $(\boldsymbol\alpha,\boldsymbol\beta)$) differs from the linear-model prediction, and $r_k$ becomes informative.

Radius update (standard thresholds $\tau_1=0.25,\ \tau_2=0.75,\ \gamma_1=0.5,\ \gamma_2=2$, with $\Delta_k$ clipped to $[\Delta_{\min},\,\Delta_{\max}]$):
\begin{equation}
\Delta_{k+1} \;=\;\min\!\Big(\Delta_{\max},\;\max\!\big(\Delta_{\min},\;\mathrm{rule}(r_k,\Delta_k)\big)\Big),\quad
\mathrm{rule}(r_k,\Delta_k) \;=\; \begin{cases}\gamma_2\Delta_k & r_k \ge \tau_2\\ \Delta_k & \tau_1 \le r_k < \tau_2\\ \gamma_1\Delta_k & r_k < \tau_1\end{cases}
\end{equation}

\paragraph{Step acceptance.} Accept $C^{(k+1)} = C^{(k)} + d_k^\star$ if $r_k \ge \tau_1$; otherwise $C^{(k+1)} = C^{(k)}$ (Yuan 2015, \S 2 step-acceptance rule).

\subsection{Bound handling via Coleman--Li affine scaling (Yuan 2015, \S 3)}

To respect $C \ge 0$ in a principled way (rather than post-hoc clipping), use the Coleman--Li scaling matrix
\[
D_k \;=\; \mathrm{diag}\!\left(\sqrt{|v^{(i)}(C^{(k)})|}\right),\qquad
v^{(i)}(C) \;=\; \begin{cases} C_i & L_i > 0 \text{ and } C_i \text{ has lower bound } 0 \\ +\infty & \text{otherwise} \end{cases}
\]
and replace the Euclidean TR with the ellipsoidal TR
\[
\tfrac12\|D_k^{-1} d\|_2^2 \;\le\; \Delta_k.
\]
This is eq.\ 3.34--3.37 of the Yuan survey. The net effect: as any component $C_i^{(k)}$ approaches its lower bound $0$, the effective step size in direction $i$ shrinks automatically, so the iterate cannot cross the bound. For $K = 2$ this reduces to a two-line scalar rule.

\subsection{Sanity checks}

\paragraph{Check 1 --- direction is cost-sensitive.} Because $L \ge 0$ componentwise with $L_k > 0$ for any class actually incurring loss, $d_k^{\mathrm{lin}} = \sqrt{2\Delta_k}\,L/\|L\|_2 \ge 0$. Every component of $C$ increases on a TR step, with the class that currently contributes the most loss gaining the most weight. \checkmark

\paragraph{Check 2 --- $\Delta_{\max} \to 0$ recovers fixed cost-sensitive learning.} If $\Delta_{\max} \to 0$, the outer ball pins $C$ to $C^0$ and every $\Delta_k \le \Delta_{\max}$ also vanishes, so $d_k^\star \to 0$. $C$ is frozen at the inverse-frequency baseline. \checkmark

\paragraph{Check 3 --- step on the boundary.} When the bound and outer-ball constraints are inactive, the linear maximizer lives on the boundary of the $\Delta_k$-ball: $\tfrac12\|d_k^\star\|_2^2 = \Delta_k$. The working radius is fully spent each step. \checkmark

\subsection{Why this is a real iterative TR, not a one-shot projection}

When (and only when) the TR $C$-block is \emph{interleaved} with re-solves of Block A --- i.e., whenever $C$ moves, $(\boldsymbol\alpha,\boldsymbol\beta)$ is updated and so is $L$ --- the sequence $\{C^{(k)}\}$ is a genuine TR trajectory on the reduced function
\[
F(C) \;=\; \min_{\boldsymbol\alpha,\boldsymbol\beta}\;\Phi(\boldsymbol\alpha,\boldsymbol\beta,C) \;+\; \lambda \mathbf{R}_{\boldsymbol\beta}(\boldsymbol\beta).
\]
By the envelope theorem, $\nabla F(C) = L(\boldsymbol\alpha^\star(C),\boldsymbol\beta^\star(C))$ --- exactly the $L_k$ we used as the linear-model gradient. The actual reduction $F(C^{(k)}+d) - F(C^{(k)})$ is \emph{not} linear in $d$ (because $(\boldsymbol\alpha^\star,\boldsymbol\beta^\star)$ moves with $C$), so $r_k \neq 1$ in general and the radius-update rule does real work: it shrinks $\Delta_k$ whenever re-solving Block A shows the linear model over-predicted the reduction in $F$, and expands $\Delta_k$ when the model is consistently accurate. Algorithm 3 below puts this interleaving into pseudocode.

\paragraph{Remark (Option TR-Q).} For a quadratic model $m_k(d) = L_k^\top d + \tfrac12 d^\top B_k d$ we need a Hessian surrogate $B_k$. The natural choice is an L-BFGS update on the sequence $\{(C^{(k)}, L_k)\}$, approximating $\nabla^2 F(C)$. This gives superlinear local convergence but requires storing gradient/iterate pairs. For $K = 2$ the linear model is almost always adequate; we leave TR-Q as future work.

\paragraph{Remark (ARC).} Cubic regularization (Yuan 2015, \S 7, eq.\ 7.2) replaces the ball constraint by a $\tfrac13\sigma_k\|d\|_2^3$ penalty in the objective. It has better worst-case complexity than TR but needs a Hessian just like TR-Q. Noted here as an alternative; not used.

%======================================================================
\section{Algorithm 3 --- full pseudocode and numerical walkthrough}
%======================================================================

\subsection{Pseudocode (matches the style of Algorithm 2 in the paper)}

\begin{algorithm}[H]
\caption{Cost-Sensitive iSFS via BCD with Trust-Region $C$-update}
\begin{algorithmic}[1]
\Require data $\{\mathbf{X}_m, \mathbf{y}_m \text{ for } m \in \mathbf{pf}\}$, baseline $C^0$, problem radius $\Delta_{\max}$, initial TR radius $\Delta_0 \in (0,\Delta_{\max}]$, lower radius $\Delta_{\min}$, TR thresholds $\tau_1=0.25,\ \tau_2=0.75,\ \gamma_1=0.5,\ \gamma_2=2$, $\boldsymbol{\beta}$-regularization $\lambda$, $\boldsymbol{\alpha}$-regularization $\mathbf{R}_{\boldsymbol{\alpha}}$
\Ensure $\boldsymbol{\alpha}$ (per-profile weights), $\boldsymbol{\beta}$ (shared coefficients), $C$ (class costs)
\State Initialize $\boldsymbol{\beta}^{(0)}$ by fitting each source individually (paper's Algorithm 2 init).
\State Initialize $C^{(0)} \gets C^0$, \;$\Delta \gets \Delta_0$.
\For{outer iteration $t = 1, 2, \ldots$}
    \State \textbf{Block A --- fix $C^{(t-1)}$, solve $(\boldsymbol{\alpha},\boldsymbol{\beta})$ by cost-sensitive Algorithm 2:}
    \Repeat
        \For{each profile $m$} \State $\boldsymbol{\alpha}_m \gets$ FISTA on eq (16$^C$) with $C_{y_{m,j}}$ weights. \EndFor
        \State $\boldsymbol{\beta} \gets$ FISTA on $g^C(\boldsymbol{\beta}) + \lambda \mathbf{R}_{\boldsymbol{\beta}}(\boldsymbol{\beta})$, gradient $\nabla g^C(\boldsymbol{\beta}^i)$ per \S 6.2.
    \Until{inner objective stops decreasing}
    \State Record $F_{\text{old}} \gets \Phi(\boldsymbol{\alpha}^{(t)},\boldsymbol{\beta}^{(t)},C^{(t-1)}) + \lambda\mathbf{R}_{\boldsymbol{\beta}}(\boldsymbol{\beta}^{(t)})$.
    \State \textbf{Block B --- TR step on $C$ around $C^{(t-1)}$:}
    \State Compute per-class loss vector $L$:
    \[ L_k = \frac{1}{|\mathbf{pf}|}\sum_{m}\frac{1}{2 n_m}\!\!\sum_{\substack{j:\,y_{m,j}=k}}(\hat y_{m,j} - y_{m,j})^2\quad\text{at } (\boldsymbol{\alpha}^{(t)},\boldsymbol{\beta}^{(t)}). \]
    \State Propose TR step (unconstrained Cauchy direction for a linear model):
    \[ d_{\mathrm{lin}} \gets \sqrt{2\Delta}\;L/\|L\|_2. \]
    \State Clip $d_{\mathrm{lin}}$ so that $C^{(t-1)}+d \ge 0$ and $\tfrac12\|C^{(t-1)}+d-C^0\|_2^2 \le \Delta_{\max}$ (scalar clip for $K=2$; Coleman--Li scaling per \S 7.4 for general $K$). Call the clipped step $d^\star$.
    \State Candidate: $\widehat C \gets C^{(t-1)} + d^\star$.
    \State Evaluate actual reduction by briefly re-solving Block A at $\widehat C$:\;run a few FISTA sweeps of cost-sensitive Algorithm 2 with $C = \widehat C$ to get $(\widehat{\boldsymbol\alpha},\widehat{\boldsymbol\beta})$, and set $F_{\text{new}} \gets \Phi(\widehat{\boldsymbol\alpha},\widehat{\boldsymbol\beta},\widehat C) + \lambda\mathbf{R}_{\boldsymbol{\beta}}(\widehat{\boldsymbol\beta})$.
    \State Predicted change (linear model on $\Phi_C$): $m(d^\star) - m(0) \gets L^\top d^\star$.
    \State $r \gets (F_{\text{new}} - F_{\text{old}}) \,/\, (L^\top d^\star)$. \Comment{sign convention: we are \emph{increasing} $\Phi_C$ in the $C$-block}
    \If{$r \ge \tau_1$} \State Accept: $C^{(t)} \gets \widehat C$. \Else \State Reject: $C^{(t)} \gets C^{(t-1)}$. \EndIf
    \State Update TR radius: $\Delta \gets \min\big(\Delta_{\max},\,\max(\Delta_{\min},\,\mathrm{rule}(r,\Delta))\big)$, where
    \[ \mathrm{rule}(r,\Delta) = \begin{cases}\gamma_2\Delta & r \ge \tau_2 \\ \Delta & \tau_1 \le r < \tau_2 \\ \gamma_1\Delta & r < \tau_1\end{cases} \]
    \If{$\|C^{(t)} - C^{(t-1)}\|_\infty < \epsilon$ \textbf{and} $\|\boldsymbol{\beta}^{(t)} - \boldsymbol{\beta}^{(t-1)}\|_2 < \epsilon$ \textbf{and} $\Delta \le \Delta_{\min}$}
        \State \Return $\boldsymbol{\alpha}^{(t)},\boldsymbol{\beta}^{(t)},C^{(t)}$.
    \EndIf
\EndFor
\end{algorithmic}
\end{algorithm}

Two outer blocks, with the $C$-block implemented as a TR iteration. The TR radius $\Delta$ adapts automatically: if a step of size $\sqrt{2\Delta}$ in direction $L/\|L\|_2$ actually reduces the saddle objective as much as the linear model predicted, $\Delta$ grows; if re-solving Block A reveals the linear model over-predicted reduction, $\Delta$ shrinks; if the reduction is too small or negative, the step is rejected.

\paragraph{Numerical walkthrough.} Deferred to a follow-up once a TR-L reference implementation exists. A worked trajectory for a synthetic 3-profile dataset will report, per outer iteration: $(\boldsymbol\alpha,\boldsymbol\beta)$ after Block A, per-class losses $L$, proposed step $d^\star$, reduction ratio $r_k$, accept/reject decision, and the updated radius $\Delta$.

%======================================================================
\section{Connection to BCD theory}
%======================================================================

\subsection{Multi-convex structure}

Cost-sensitive iSFS sits inside the \textbf{multi-convex} family from \texttt{BCD\_explained.md} \S 3 / Aggarwal \S 4.10.2:

\begin{center}
\begin{tabular}{lllp{0.40\linewidth}}
\toprule
\textbf{block fixed} & \textbf{block updated} & \textbf{sub-problem} & \textbf{reason convex (or concave)} \\
\midrule
$\boldsymbol{\beta}, C$ & $\boldsymbol{\alpha}_m$ per profile & per-profile weighted least-squares & positive-semidefinite quadratic in $\boldsymbol{\alpha}_m$ \\
$\boldsymbol{\alpha}, C$ & $\boldsymbol{\beta}$ & weighted regularized least-squares & positive-definite quadratic in $\boldsymbol{\beta}$ \\
$\boldsymbol{\alpha}, \boldsymbol{\beta}$ & $C \in \mathcal B(C^0,\Delta_{\max})\cap\mathbb{R}^K_{\ge 0}$ & linear-in-$C$ over a convex compact set, solved by TR iteration & linear function on a ball (plus bounds) \\
\bottomrule
\end{tabular}
\end{center}

The $(\boldsymbol\alpha,\boldsymbol\beta)$ blocks have closed-form-feeling solves (FISTA); the $C$-block is a genuine TR iteration whose \emph{inner step} is closed-form (the linear Cauchy step) but whose outer radius adaptation makes the overall $C$-update iterative. No learning rate inside any block.

\subsection{What converges and what doesn't}

Two facts coexist:

\begin{enumerate}[leftmargin=*,itemsep=2pt]
  \item \textbf{Inner BCD on $(\boldsymbol{\alpha}, \boldsymbol{\beta})$ with $C$ fixed converges monotonically.} This is the standard multi-convex BCD theorem from Aggarwal \S 4.10.2: each step decreases the inner objective, so the inner loop is monotone.
  \item \textbf{The outer trajectory is \emph{not} a strict monotonic decrease in $\Phi$} --- the C-block is a $\max$, so it can increase $\Phi$ between iterations even though the $(\boldsymbol{\alpha}, \boldsymbol{\beta})$ block decreases it. The right invariant for the outer loop is the \textbf{saddle-point gap}:
\end{enumerate}
\[
\text{gap}(\boldsymbol{\alpha}, \boldsymbol{\beta}, C) \;=\; \max_{C' \in \mathcal B(C^0,\Delta_{\max}),\,C'\ge 0}\Phi(\boldsymbol{\alpha},\boldsymbol{\beta},C') \;-\; \min_{\boldsymbol{\alpha}',\boldsymbol{\beta}'}\Phi(\boldsymbol{\alpha}',\boldsymbol{\beta}',C)
\]

Under standard conditions (strong convexity in $(\boldsymbol{\alpha}, \boldsymbol{\beta})$ and the bounded-set constraint on $C$), saddle-point alternation reduces the gap to zero. Moreover, the TR radius update (Yuan 2015, \S 2) guarantees that the $C$-block behaves like a descent method on the reduced function $F(C) = \min_{\alpha,\beta} \Phi$: the radius shrinks whenever the linear model over-predicts, so the algorithm cannot oscillate indefinitely.

\subsection{Connection to the textbook BCD case studies}

\begin{itemize}[leftmargin=*,itemsep=2pt]
  \item \textbf{K-means as BCD} (Aggarwal \S 4.10.3): two blocks, centroids and assignments, each with a closed-form optimum (mean and nearest-prototype). Cost-sensitive iSFS does the same with $(\boldsymbol{\alpha}, \boldsymbol{\beta})$ and $C$, using weighted normal equations for the $(\boldsymbol\alpha,\boldsymbol\beta)$-block and a trust-region iteration for the $C$-block.
  \item \textbf{ALS as BCD} (Aggarwal \S 8.3.2.3): two blocks $(U, V)$, each row solved as an independent ridge regression. Cost-sensitive iSFS's $\boldsymbol{\beta}$ update is structurally the same --- one global weighted ridge solve.
  \item \textbf{\texttt{opt.impute}} (Bertsimas et al.\ JMLR 2018): two blocks, auxiliary state $U$ and imputations $(W, V)$, alternated. Cost-sensitive iSFS adds a third ``auxiliary state'' ($C$) and treats it the same way.
\end{itemize}

The trade-off Aggarwal flags (\S 4.10.2) applies directly:

\begin{quote}
``Each step in block coordinate descent is more expensive, but fewer steps are required.''
\end{quote}

For cost-sensitive iSFS at moderate $n$ and small $K$, each block is cheap: $(\boldsymbol\alpha,\boldsymbol\beta)$ via FISTA, $C$ via a closed-form Cauchy step plus scalar radius update.

%======================================================================
\section{Quick reference card}
%======================================================================

\begin{center}
\small
\begin{tabular}{lll}
\toprule
\textbf{symbol} & \textbf{meaning} & \textbf{where it lives} \\
\midrule
$\alpha^i_m \in \mathbb{R}$ & source-$i$ weight for profile $m$ & inner block, $\mathbf{R}_{\boldsymbol{\alpha}}(\boldsymbol{\alpha}_m) \leq 1$ \\
$\boldsymbol{\beta}^i \in \mathbb{R}^{p_i}$ & shared feature coefficients for source $i$ & inner block, $\mathbf{R}_{\boldsymbol{\beta}}$ regularized \\
$C \in \mathbb{R}^K_{\ge 0}$ & per-class cost vector & outer block, $C \in \mathcal B(C^0,\Delta_{\max})$ \\
$C^0 \in \mathbb{R}^K_{\ge 0}$ & rescaled inverse-frequency baseline & user-chosen prior, TR center \\
$\Delta_{\max}$ & problem-level TR radius around $C^0$ & hyperparameter (cross-validated) \\
$\Delta_k$ & working TR radius at step $k$ of the $C$-block & updated by Yuan ratio rule \\
$r_k$ & actual/predicted reduction ratio & governs $\Delta_k$ update \\
$\lambda$ & strength of $\mathbf{R}_{\boldsymbol{\beta}}$ & hyperparameter (cross-validated) \\
$L_k$ & per-class loss coefficient (gradient of $\Phi_C$ in class $k$) & recomputed every $C$ step \\
$\hat{y}_{m,j}$ & per-sample model prediction $\sum_i \alpha^i_m \mathbf{x}^i_{m,j} \cdot \boldsymbol{\beta}^i$ & computed inside both blocks \\
$C_{y_{m,j}}$ & class-cost lookup for sample $j$ in profile $m$ & scalar from the $K$-vector $C$ \\
\bottomrule
\end{tabular}
\end{center}

\bigskip

\begin{center}
\small
\begin{tabular}{lp{0.50\linewidth}p{0.30\linewidth}}
\toprule
\textbf{update} & \textbf{formula} & \textbf{meaning} \\
\midrule
$\boldsymbol{\alpha}_m$ & $\arg\min$ of $\frac{1}{2}\sum_j C_{y_{m,j}}(\hat{y}-y)^2$ s.t.\ $\mathbf{R}_{\boldsymbol{\alpha}} \leq 1$ & weighted per-profile constrained LS \\
$\boldsymbol{\beta}$ & $\arg\min$ of $g^C(\boldsymbol{\beta}) + \lambda \mathbf{R}_{\boldsymbol{\beta}}(\boldsymbol{\beta})$ via FISTA, gradient $\nabla g^C(\boldsymbol{\beta}^i)$ & weighted regularized LS \\
$C$ & TR-L step: $d_{\mathrm{lin}} = \sqrt{2\Delta_k}\,L/\|L\|_2$, clipped to $C\ge 0$ and $\mathcal B(C^0,\Delta_{\max})$; accept if $r_k \ge \tau_1$ & trust-region iteration on $C$ (Yuan 2015 \S 2) \\
\bottomrule
\end{tabular}
\end{center}

\bigskip

\begin{center}
\small
\begin{tabular}{ll}
\toprule
\textbf{theoretical fact} & \textbf{source} \\
\midrule
BCD with closed-form block updates monotonically decreases the inner objective on multi-convex problems & Aggarwal \S 4.10.2; \texttt{BCD\_explained.md} \S 3 \\
DRO with an $\ell_2$ / $\chi^2$ ball around the class-distribution prior is principled cost-sensitive learning & Hu, Niu, Sato, Sugiyama (2018); Namkoong \& Duchi (2017) \\
Trust-region methods solve bounded-step subproblems with adaptive radius via the actual/predicted ratio & Yuan (2015), \S 2, eq.\ 2.3--2.7 \\
Coleman--Li affine-scaled TR handles bound constraints ($C \ge 0$) inside the TR subproblem & Yuan (2015), \S 3.3, eq.\ 3.34--3.37 \\
The original iSFS Algorithm 2 (2-block BCD on $\boldsymbol{\alpha}, \boldsymbol{\beta}$) is the inner loop here & Xiang et al.\ (2013), \S 3.2 \\
$K$-means and ALS are the canonical 2-block BCD examples this method generalizes & Aggarwal \S 4.10.3 and \S 8.3.2.3; \texttt{BCD\_explained.md} \S 4--5 \\
\bottomrule
\end{tabular}
\end{center}

\bigskip

\begin{center}
\emph{Cost-sensitive iSFS with a trust-region $C$-block is a strict generalization of the original iSFS model: setting $\Delta_{\max} \to 0$ pins $C$ to the inverse-frequency baseline and recovers classical fixed-cost cost-sensitive learning, while finite $\Delta_{\max}$ lets the optimization adapt $C$ from the data via an iterative TR step (Yuan 2015, \S 2) that fits inside the same BCD machinery the rest of the model already uses.}
\end{center}

%======================================================================
\appendix
\section{Trust-region background (focused catalog)}
%======================================================================

Drawn from Yuan, Y.\ (2015), \emph{Recent advances in trust-region algorithms}, Math.\ Program.\ Ser.\ B 151, 249--281. We list the variants that bear on our $C$-block.

\subsection{A.1\; The canonical trust-region subproblem (Yuan \S 2)}

At iterate $x_k$ with gradient $g_k = \nabla f(x_k)$ and Hessian-surrogate $B_k$, the standard TR subproblem (eq.\ 2.3--2.4, p.\ 252) is
\begin{equation}
\min_{d \in \mathbb{R}^n}\; m_k(d) \;=\; f(x_k) + g_k^\top d + \tfrac12 d^\top B_k d\qquad\text{s.t.}\quad \|d\|_2 \le \Delta_k.
\tag{TR}
\end{equation}
Each step defines
\[
\text{Pred}_k = m_k(0) - m_k(d_k),\qquad \text{Ared}_k = f(x_k) - f(x_k + d_k),\qquad r_k = \frac{\text{Ared}_k}{\text{Pred}_k}
\]
(eq.\ 2.7, p.\ 253). The step is accepted if $r_k \ge \tau_0$ (commonly $\tau_0 = \tau_1$), and the radius is updated by
\[
\Delta_{k+1} = \begin{cases}\gamma_2 \Delta_k & r_k \ge \tau_2 \\ \Delta_k & \tau_1 \le r_k < \tau_2 \\ \gamma_1 \Delta_k & r_k < \tau_1\end{cases}
\]
with typical $(\tau_1,\tau_2,\gamma_1,\gamma_2) = (0.25,\,0.75,\,0.5,\,2)$. The \textbf{Cauchy decrease} (sufficient-descent) condition (eq.\ 2.5, p.\ 252) is
\[
\text{Pred}_k \;\ge\; c\,\|g_k\|_2\,\min\!\Big\{\Delta_k,\,\|g_k\|_2/\|B_k\|_2\Big\},\qquad c > 0.
\]

\textbf{How we use this.} See the mapping table in \S 5.8 and the convention note in \S 5.9 (the main body). The $C$-block uses $B_k = 0$ (linear model) with our squared-distance convention $\tfrac12\|d\|_2^2 \le \Delta_k$.

\subsection{A.2\; Affine-scaled TR for bound constraints (Yuan \S 3.3, eq.\ 3.34--3.37)}

For a bound-constrained problem $\min f(x)$ s.t.\ $\ell \le x \le u$, Coleman \& Li (1996) replace the Euclidean ball with an ellipsoid whose axes shrink in directions where the current iterate is near a bound. Define
\[
v^{(i)}(x) \;=\; \begin{cases} x_i - \ell_i & \text{if } (\nabla f(x))_i \ge 0 \text{ and } \ell_i > -\infty \\ u_i - x_i & \text{if } (\nabla f(x))_i < 0 \text{ and } u_i < +\infty \\ +\infty & \text{otherwise} \end{cases}
\]
and set the scaling $D_k = \mathrm{diag}\big(\sqrt{|v^{(i)}(x_k)|}\big)$. The TR subproblem becomes
\[
\min\; m_k(d)\qquad\text{s.t.}\quad \tfrac12\|D_k^{-1} d\|_2^2 \le \Delta_k.
\]
In directions where $v^{(i)} \to 0$ (a bound is active), the effective step shrinks to zero, so no iterate can cross a bound. \textbf{How we use this.} For $C \ge 0$ with the gradient $L \ge 0$ (so the ascent direction never points \emph{toward} zero except on components where $L_i = 0$), scaling is usually inactive and the Euclidean step works. The Coleman--Li scaling is the principled fallback when any $C_i^{(k)}$ does approach $0$.

\subsection{A.3\; Other variants --- noted but not used here}

\begin{itemize}[leftmargin=*,itemsep=2pt]
  \item \textbf{Retrospective TR radius rule} (Bastin et al., eq.\ 2.13--2.14, p.\ 255). Updates $\Delta_{k+1}$ using the \emph{next} model $m_{k+1}$ rather than $m_k$. Useful when model quality changes rapidly between iterates. \emph{Not used}: our linear model is constant within a $C$-block, so retrospective and prospective updates coincide.
  \item \textbf{Adaptive-radius TR} $\Delta_k = \mu_k\|g_k\|_2$ (Fan \& Yuan, eq.\ 2.12, p.\ 254). Scales the radius with current gradient magnitude. \emph{Noted}; we instead start $\Delta_0$ as a hyperparameter.
  \item \textbf{Non-monotone TR} (Toint; Yuan \S 2). Compares reduction to the max of the last $M$ iterates rather than the immediate predecessor. \emph{Relevant for saddle problems}: the outer BCD objective need not decrease monotonically. Noted as a possible future robustness improvement.
  \item \textbf{Subspace TR} (Yuan eq.\ 2.16--2.21, p.\ 256). Restricts $d$ to a low-dimensional subspace (span of recent gradients/steps). \emph{Not used}: $K = 2$ is already 2-dimensional.
  \item \textbf{Conic-model TR} (Di \& Sun; Yuan eq.\ 2.6, p.\ 252). Replaces the quadratic model with a rational/conic one to preserve nonconvex structure. \emph{Not used}: the $C$-block is already linear.
  \item \textbf{Cubic regularization (ARC)} (Cartis--Gould--Toint; Yuan eq.\ 7.2, p.\ 272). Replaces the ball constraint with a $\tfrac13 \sigma_k\|d\|^3$ penalty. Better worst-case iteration complexity than TR. \emph{Noted as an alternative to the quadratic-model variant TR-Q in \S 7 above}; requires a Hessian surrogate we do not currently maintain.
\end{itemize}

\subsection{A.4\; Subproblem solvers (not needed for $K = 2$, cited for completeness)}

\begin{itemize}[leftmargin=*,itemsep=2pt]
  \item \textbf{Steihaug truncated-CG} (Steihaug 1983; Yuan p.\ 252). Runs CG on $(B_k,-g_k)$ and terminates at the first of three conditions: CG convergence, direction of negative curvature detected, or iterate hits the ball boundary. Achieves at least a Cauchy-fraction of the optimal reduction.
  \item \textbf{Mor\'e--Sorensen} (Mor\'e \& Sorensen 1983). Solves the secular equation $\|(B_k + \lambda I)^{-1}g_k\|_2 = \Delta_k$ for the Lagrange multiplier $\lambda^\star$, giving the globally optimal TR step. Globally convergent but more expensive.
  \item \textbf{Dogleg} (Powell 1970). Two-piece path combining steepest descent and Newton directions; exact for convex quadratics with $B_k \succ 0$.
\end{itemize}

For our TR-L subproblem on $K = 2$ the solver is a single scalar: $d^\star = \sqrt{2\Delta_k}\,L/\|L\|_2$ followed by a one-line clip to $C \ge 0$.

\end{document}
